\subsection{Microdosing LSD}
\totalscore{8}

Microdosing psychedelic drugs like LSD recently gained increasing popularity.
Microdosing is the practice of regularly using low doses of such psychedelic drugs, e.g.\ $\sim5\%$-$20\%$ of a typical recreational dose of LSD for about
one to three times per week. Anecdotal reports claim that microdosing improves well-being, creativity and cognition. 
During this exercise we want to have a causal look at two recent studies\footnote{For the purpose of this exam results are (over-)simplified and possibly misrepresented. We do not guarantee for the correctness of the representation of the studies in this exam. References are below in the footnotes.} that investigated such effects.

\subsubsection*{Study A\footnote{R.\ Petranker et al., \textit{Microdosing psychedelics: Subjective benefits and challenges, substance testing behavior, and the relevance of intention}, Journal of Psychopharmacology, Oct.\ 2020. https://doi.org/10.1177/0269881120953994.}}

The researchers of study A wrote:
\begin{quote}
``Method: The Global Drug Survey runs the world's largest online drug survey. Participants who reported last 
year use of [LSD] in the Global Drug Survey 2019 were offered the opportunity 
to answer a sub-section on microdosing. 
  [..] 
  We asked participants to indicate which benefits from the following 
list applied to their experience of microdosing: (a) enhanced mood, (b) enhanced focus, (c) enhanced creativity, [..] (f) enhanced sociability [etc.].

  [..]

Results: Data from [4769] people who reported microdosing at least once in the last 12 months were used for analyses.
Our results suggest a [..] replication of previously reported benefits [..:] enhanced mood, creativity, focus and sociability.''
\end{quote}
%``Conclusion: Our results suggest the perceived benefits associated with microdosing greatly outweigh the challenges. Microdosing may have utility for a variety of uses while having minimal side effects.''

\begin{table}[h!]
\caption{Study A: Reported benefits for microdosing LSD. Total number of responses: 4769.
\label{tab:study-A}}
\begin{center}
\begin{tabular}{ l|l|l }
    Reported Benefit & Count & Percentage \\
 \hline
   Enhanced mood & 2443 & 51.23 \\
   Encanced creativity & 2223 & 46.61\\
   Enhanced sociability & 1797 & 37.68\\
   Enhanced focus & 1561 & 32.73
\end{tabular}
\end{center}
\end{table}

Carefully read the above summary and look at the table in Table~\ref{tab:study-A}.
For simplicity, you may only focus on the outcome variable `enhanced mood' in answering the next question.
\begin{enumerate}[label=\alph*)]
  \item Explain what quantity one would like to estimate, what quantity the researchers actually measured, and how they came to their conclusion. 
Explain the found association between the relevant variables. % from several points of views. 
Discuss causal relations between the relevant variables. 
You may draw a causal Bayesian network to aid your explanation.
Discuss shortcomings in the design and evaluation of this study from a causal point of view. 
\score{4}
\answer{
    Let $Y$ be the outcome variable (e.g.\ enhanced mood),
     $X$ be the amount (or indicator) of microdosing in the past year,
     $T$ the indicator variable if the person participated in the Global Drug Survey,
     $S$ the indicator variable if the person participated in the microdosing study.
     One is interested in the causal effect: $P(Y|\doit(X))$, i.e.\ in the difference between $P(Y|\doit(X=0))$ and $P(Y|\doit(X=1))$.
     The researchers estimated the distribution of $P(Y|T=1,S=1,X=1)$, where they modelled $Y$ as binary.

    A possible CBN could be:
    \begin{center}\begin{tikzpicture}
      \node[var] (Y) at (2,0) {$Y$};
      \node[var] (X) at (0,0) {$X$};
      \node[var] (S) at (1,-1.5) {$S$};
      \node[var] (T) at (-1,-1.5) {$T$};
      \draw[arr] (T) to (S);
      \draw[arr] (X) to (Y);
      \draw[arr] (X) to (S);
      \draw[arr] (X) to (T);
      \draw[biarr,bend left] (X) edge (Y);
    \end{tikzpicture}\end{center}

    This is solely an observational study. Furthermore, the data underlies at least two types of selection bias (only people in the drug survey, $T=1$, and only people who already microdose, indicated by $S=1$). Also no comparison group ($X=0$) was used.
There could be further counfounding variables, e.g.\ possibly only persons who take drugs experience enhanced mood due to countering withdrawal effects, etc.
}
\points{4 points. It is hard to be specific at this stage about how to assign points to partially correct/incomplete answers.
    What we want to measure: 1 pt for measuring causal effect $P(Y|\doit(X=1))-(Y|\doit(X=0))$ as goal (0.5 pt for $P(Y|\doit(X))$ and none for $P(Y|X=1)$).\\
    What they measured: 1 pt for $P(Y|X=1,S=1,T=1)$  (0.5pt if some of the variables $S,T$ on the conditioning side are missing).\\
    Reasonal explanation or drawing of causal Bayesian network: 1pt.\\
    Explicitely pointing out that the study is non-interventional/solely observational: 0.5 pt.
    Pointing out that there was not even a comparison group: 0.5 pt.\\
After adding all points together, round up to integer points.
}
\end{enumerate}

\subsubsection*{Study B\footnote{B.\ Szigeti et al., \textit{Self-blinding citizen science to explore psychedelic microdosing}, eLife, Mar.\ 2021. https://elifesciences.org/articles/62878.}}

The researchers of study B designed their study as follows.

LSD users were recruited through advertisement on relevant online and offline forums. 
A total of 191 participants completed the study.
Participants were randomly assigned to either a microdosing group (MD) or a placebo group (PL) with equal probability.
During 4 weeks each participant was given two pills per week, each either containing a microdose of LSD if in group MD or a
placebo if in group PL. 
Before, during and after these weeks the participants were tested for cognitive performance, mental well-being,
social connectedness, life satisfaction, etc.
Furthermore, each participant was regularly asked to guess if they were in the MD or the PL group.

\begin{table}[h!] %this is figure 5 from the study. numbers are just read off the figure there and rounded and rescaled
\caption{Study B: Mean and 95\% confidence interval of the corresponding score (higher is better) 
    adjusted for actual group and the guess of the participant in which group they are. 
    The corresponding p-values for a mean difference test are either below 0.01 (very low) if the 
  95\%-confidence intervals have no overlaps and between 0.1-0.9 (high) if they have overlap.\label{tab:study-B}}
\begin{center}
\begin{tabular}{ c|c|l|l|l|l}
    Group & Guess  & Emotion Score & Mood Score & Well-Being Score & Cognitive Score \\
 \hline
    PL & PL & 12.8 $\pm$ 1.1 & 57 $\pm$ 2& 47.9 $\pm$ 1.1 & -9 $\pm$ 7\\
    MD & PL & 13.9 $\pm$ 1.2 & 58 $\pm$ 3& 48.1 $\pm$ 1.2 & -8 $\pm$ 9\\
    PL & MD &  18.1 $\pm$ 1.9& 69 $\pm$ 4& 51.1 $\pm$ 1.2 & -10 $\pm$ 9\\
    MD & MD &   19.1 $\pm$ 1.4&72 $\pm$ 3& 50.9 $\pm$ 1.0 & -1 $\pm$ 9
\end{tabular}
\end{center}
\end{table}

Have a look at Table~\ref{tab:study-B}. It shows an extract of the estimated mean and 95\%-confidence intervals
of each score (emotional, mood, well-being, cognitive performance, etc.) adjusted for the actual group the participants were
in and also for the group the participants guessed they were in.
For simplicity, you may only focus on the outcome variable `mood' in answering the next question.

\begin{enumerate}[resume*]
  \item 
    Answer the same questions as in a) but now for study B.
    \score{4}
\answer{
    Let $Y$ be the outcome variables (mood score).
    We want to estimate the difference between $E(Y|\doit(X=1))$ and $E(Y|\doit(X=0))$ ($E$ denoting expectation).
    The researchers of study B estimated $E(Y|G,T=1,\doit(X))$, where $G$ is the guess of the participant.
    There is (pre-intervention) selection bias $T=1$ (as only LSD users were recruted).
    A possible SCM could have the following graph:
    \begin{center}\begin{tikzpicture}
      \node[var] (Y) at (2,0) {$Y$};
      \node[varc] (X) at (0,0) {$X$};
      \node[var] (G) at (1,-1.5) {$G$};
      \node[var] (T) at (1,1.5) {$T$};
      \draw[arr] (X) to (Y);
      \draw[arr] (T) to (Y);
      \draw[arr,bend left] (G) edge (Y);
      \draw[arr,bend left] (Y) edge (G);
      \draw[biarr] (Y) edge (G);
    \end{tikzpicture}\end{center}
    The study is an interventional study, more specifically, a blind randomized controlled trial (the provided text doesn't say if it is double blind or not).
    The results can be interpreted as that the expectation the participant has (guess $G=1$) is associated with outcome,
    but not the actual microdose itself. Cognitive performance is unaltered. 
    From the numbers alone we could not conclude that $G$ causes $Y$ as we only condition on $G$.
    Whether we can restrict the possible causal relations between $G$ and $Y$ is not clear from the description of the experimental procedure (is $G$ measured before $Y$, or after $Y$?).
    %but we could allow background knowledge about that expectations can causally change well-being, etc., to come to this conclusion.
}
\points{4 points. It is hard to be specific at this stage about how to assign points to partially correct/incomplete answers.\\
What we want to measure: 1 pt for measuring causal effect $P(Y|\doit(X=1))-(Y|\doit(X=0))$ as goal (0.5 pt for $P(Y|\doit(X))$ and none for $P(Y|X=1)$).\\
What they did measure: 1 pt for $P(Y|G,S=1,\doit(X))$ (table 2) (0.5 pt if some variables are missing on conditioning side).\\
Explaining the association in table 2 in words: 1 pt.\\
Explicitely pointing out that we have an interventional study: 0.5 pt,\\
discussing that we still might have selection bias and not every quantity is provided to estimate the causal effect: 0.5 pt.\\
After adding all points up, rounding up to interger points.
}
\end{enumerate}
